% !TeX program = xelatex
\documentclass[UTF8,fontset=fandol,12pt,a4paper]{ctexart}
\usepackage[margin=2.5cm]{geometry}
\usepackage{amsmath,amssymb,booktabs,array,longtable,graphicx,float,caption}
\usepackage{enumitem,fancyvrb,fvextra,xcolor,placeins}
\usepackage{needspace}
\usepackage[hidelinks]{hyperref}
\IfFontExistsTF{SimSun}{\setCJKmainfont[AutoFakeBold]{SimSun}\setCJKsansfont[AutoFakeBold]{SimHei}}{}
\IfFontExistsTF{Times New Roman}{\setmainfont{Times New Roman}}{}
\IfFontExistsTF{Consolas}{\setmonofont{Consolas}}{}
\linespread{1.22}
\setlength{\parindent}{2em}
\setlength{\parskip}{0.15em}
\setlength{\textfloatsep}{12pt plus 2pt minus 2pt}
\setlength{\intextsep}{10pt plus 2pt minus 2pt}
\setlength{\abovedisplayskip}{8pt plus 2pt minus 2pt}
\setlength{\belowdisplayskip}{8pt plus 2pt minus 2pt}
\ctexset{section={format=\large\bfseries,number=\chinese{section},beforeskip=16pt,afterskip=8pt},subsection={format=\normalsize\bfseries,beforeskip=10pt,afterskip=5pt}}
\setlist[enumerate]{nosep,leftmargin=2em}
\captionsetup{font=small,labelfont=bf,labelsep=quad}
\graphicspath{{figures/}}
\pagestyle{plain}
\raggedbottom
\renewcommand{\topfraction}{0.9}
\renewcommand{\bottomfraction}{0.85}
\renewcommand{\textfraction}{0.08}
\renewcommand{\floatpagefraction}{0.75}
\makeatletter
\setlength{\@fptop}{0pt}
\setlength{\@fpsep}{16pt}
\setlength{\@fpbot}{0pt plus 1fil}
\makeatother
\newcommand{\dd}{\,\mathrm d}
\newcommand{\unit}[1]{\,\mathrm{#1}}
\newcommand{\figone}[4]{\begin{figure}[htbp]\centering\includegraphics[width=#2\linewidth]{#1}\caption{#3}\label{#4}\end{figure}}
\newcommand{\figpair}[4]{\begin{figure}[htbp]\centering\includegraphics[width=.49\linewidth]{#1}\hfill\includegraphics[width=.49\linewidth]{#2}\caption{#3}\label{#4}\end{figure}}
\hypersetup{pdftitle={考虑变物性与径向收缩的药材热湿传递及烘干时长预测},pdfauthor={}}

\hypersetup{pdftitle={AI工具使用详情},pdfauthor={}}
\begin{document}
\small\linespread{1.08}\selectfont
\begin{center}{\Large\bfseries AI工具使用详情}\end{center}
\section*{一、可确认的工具与使用范围}
2026年9月12日，本次论文整理在OpenAI Codex中进行，会话系统标识为GPT-6。AI读取工作区内既有题目、模型说明、公式推导、Python代码、计算结果和验证报告，辅助组织中文正文、检查公式与程序口径、生成LaTeX排版以及从未舍入数组生成题定表格。另使用独立上下文核查流程，通过出版社、作者机构和SciPy官方资料复核五条参考文献。

本次任务没有重新设计四问物理模型，没有将未完成的创新建议写成实验结论，也没有在起草时重新运行四问的完整PDE生产计算。全文采用既有q1-closeout-v1、q2-closeout-v1、q3-closeout-v1和q4-baseline-v1结果。具体数值核验的执行日期与范围以随附验证记录为准。

\section*{二、主要提示方式与交互过程}
用户首先要求查看项目文件夹，随后提出：“请你整体查看一下我目前的工作，根据我目前的模型建立和公式的推到写一篇条理清晰符合国赛标准的论文”。本次据此恢复此前暂停的写作，要求正文建立在现有模型和结果上。

主要处理流程为：清点题面与已有材料；对照现行代码辨别历史方案；核对输入插值、物性函数、有限体积与材料坐标方程；根据未舍入源数组生成六张题定表；撰写各问推导和结果解释；编译PDF并检查公式、图表、引用、页数与匿名性。文献核查仅将已取得的书目、摘要或官方算法说明用于相应强度的论述。

\section*{三、AI输出的采用与核验记录}
本次采用AI生成的章节组织、中文行文和LaTeX源码；表格数值由程序读取既有NPZ数组后格式化，避免人工转抄。正文修正了历史材料中的插值和网格描述不一致，保留PCHIP及实际表面加密网格；第三问使用57.4731 h、第四问使用51.0877 h，并解释末行含水率0.1500的舍入含义。未把两问时长差全部归因于收缩，未把数值收敛表述为真实药材实测验证。

自动核验包括：题定表与源数组逐项一致性、正文数值与当前验证报告对照、公式和交叉引用检查、PDF编译与页面视觉检查、程序附录的语法检查及文件一致性记录。自动核验不等同于参赛队员已经逐条完成的人工核验。本次对话尚未收到参赛队员确认最终全文已完成逐项人工审阅的记录，因此本说明不代作此项声明。

\section*{四、历史使用记录的边界}
工作区包含此前AI协作记录，以及以“Claude评价”命名的外部审阅文档，但现有材料不足以完整确认所有历史工具的具体版本、提示过程与人工采用情况。本文件准确记录本次可确认的使用范围。参赛队提交前应依据实际历史记录补充此前工具、用途和人工修改核验情况，并核对本说明与最终论文、支撑材料的一致性。

\section*{五、说明依据}
本说明依据2026年国赛AI工具使用规定的公开信息组织，论文参考文献前已列AI工具使用声明。本文件不含队员、学校或赛区身份信息。
\noindent\url{https://dxs.moe.gov.cn/zx/a/hd_sxjm_gsyw/260807/2051819.shtml}
\end{document}
